% 文档使用中文 article 类、A4 纸和 11pt 正文字号；字体在下方单独指定。
\documentclass[UTF8,a4paper,11pt,fontset=none]{ctexart}
% 页面四边留白为 25 mm，页眉高度为 15 pt。
\usepackage[margin=25mm,headheight=15pt]{geometry}
% 正文西文使用 Times New Roman，中文使用宋体（AR PL UMing CN）。
% 代码块和行内代码使用等宽西文字体，中文保持宋体。使用 XeLaTeX 编译。
\usepackage{fontspec}
\setmainfont{times new roman}[
  Path=/home/wumej22/.local/share/fonts/,
  UprightFont={times new roman.ttf},
  BoldFont={times new roman bold.ttf},
  ItalicFont={times new roman italic.ttf},
  BoldItalicFont={times new roman bold italic.ttf}
]
\setmonofont{DejaVu Sans Mono}[Scale=0.82]
\setCJKmainfont{AR PL UMing CN}[AutoFakeBold=2]
\setCJKsansfont{AR PL UMing CN}[AutoFakeBold=2]
\setCJKmonofont{AR PL UMing CN}[AutoFakeBold=2]
% 表格、颜色、长度计算及环境钩子所需宏包。
\usepackage{xcolor,booktabs,longtable,array,calc,etoolbox,upquote}
% 代码块使用浅灰蓝背景和等宽字体；长行允许自动折行。
% snugshade 提供背景，CPSCode 控制代码文字与行距。
\usepackage{fvextra,framed}
\definecolor{shadecolor}{HTML}{F2F5F8}
\DefineVerbatimEnvironment{CPSCode}{Verbatim}{fontsize=\small,breaklines=true,breakanywhere=true,breaksymbolleft={},baselinestretch=1.08}
% 页眉页脚、列表、超链接及长网址换行所需宏包。
\usepackage{fancyhdr,enumitem,hyperref,xurl}
% 标题和链接使用统一的深蓝色，并设置 PDF 元数据。
\definecolor{CPSNavy}{HTML}{163B5C}
\hypersetup{colorlinks=true,linkcolor=CPSNavy,urlcolor=CPSNavy,pdftitle={社区区域气候模式预处理系统用户手册},pdfauthor={吾买尔江·吾布力卡斯木等}}
% 章节标题字号、颜色和间距；编号与目录均显示至三级标题。
\ctexset{section={format=\Large\bfseries\color{CPSNavy},beforeskip=2.2ex,afterskip=1ex},subsection={format=\large\bfseries\color{CPSNavy}},subsubsection={format=\normalsize\bfseries},secnumdepth=3,tocdepth=3}
% 段落间距、断行容差和列表缩进。
\setlength{\parskip}{0.35em}
\setlength{\emergencystretch}{3em}
\setlist{nosep,leftmargin=2em}
% 兼容正文转换时使用的紧凑列表及图片命令。
\newcommand{\tightlist}{\setlength{\itemsep}{0pt}\setlength{\parskip}{0pt}}
\providecommand{\pandocbounded}[1]{#1}
% 页眉左侧显示手册名称，右侧显示页码，下方添加横线。
\pagestyle{fancy}
\fancyhf{}
\fancyhead[L]{\small 社区区域气候模式预处理系统用户手册}
\fancyhead[R]{\small\thepage}
\renewcommand{\headrulewidth}{0.4pt}
% 长表格缩小字号和列间距；代码背景的内边距为 8 pt。
\AtBeginEnvironment{longtable}{\footnotesize\sloppy\setlength{\tabcolsep}{3pt}}
\setlength{\FrameSep}{8pt}
\begin{document}
% 封面：调整各段前的垂直间距可改变标题、作者和版本的位置。
\begin{titlepage}
\centering
\vspace*{35mm}
{\fontsize{27}{40}\selectfont\bfseries 社区区域气候模式预处理系统\\用户手册\par}
\vspace{12mm}{\color{CPSNavy}\rule{0.85\textwidth}{0.6pt}}\par
\vspace{8mm}{\LARGE\bfseries CRESM Preprocessing System\\User Manual\par}
\vspace{23mm}{\large 吾买尔江·吾布力卡斯木　张树磊　张焓　陈虹静\par}
\vfill{\large Version 1.2.5\\[4mm]2026-09\par}
\vspace{15mm}
\end{titlepage}
% 目录使用大写罗马页码；正文另起一页并从阿拉伯数字 1 开始。
\pagenumbering{Roman}
\tableofcontents
\clearpage
\pagenumbering{arabic}
\hypertarget{ux5f15ux8a00}{%
\section{引言}\label{ux5f15ux8a00}}

社区区域气候模式预处理系统（CRESM Preprocessing System，CPS）是为新一代国产区域地球系统模式 CRESM 开发的数据前处理系统，用于统一管理和自动化执行模式运行前所需的各类数据准备任务。CRESM 基于区域气候模式 CWRF，并进一步耦合了新一代通用陆面模式 CoLM2024，通过耦合器架构实现陆---气过程的协同模拟，对数据一致性和前处理流程的规范性提出了更高要求。

在实际应用中，CRESM 的前处理工作涉及多个独立子系统和复杂的数据依赖关系，包括 CWRF 所需的静态地理数据与侧边界驱动场准备、CoLM 所需的基础地表参数场与初始状态场构建，以及 CPL7 耦合器所需的网格映射与插值权重文件生成等步骤。这些流程通常分散在不同工具和脚本中，配置复杂、操作繁琐，且缺乏统一的流程管理机制，容易产生人为配置错误和重复劳动。

CPS 将上述步骤组织为可配置的处理流程，用户可通过配置文件选择处理环节、设置实验参数和运行环境。

\hypertarget{ux7cfbux7edfux73afux5883ux8981ux6c42}{%
\section{系统环境要求}\label{ux7cfbux7edfux73afux5883ux8981ux6c42}}

CPS 系统面向高性能科学计算环境设计，主要运行于 Linux 平台，用于支持 CRESM 模式体系的数据前处理工作。在安装和使用 CPS 之前，请确保计算环境满足以下基本要求。

\hypertarget{ux64cdux4f5cux7cfbux7edfux8981ux6c42}{%
\subsection{操作系统要求}\label{ux64cdux4f5cux7cfbux7edfux8981ux6c42}}

推荐使用基于 Linux 的操作系统环境：

\begin{itemize}
\item
  Linux

  \begin{itemize}
  \tightlist
  \item
    Ubuntu 20.04 及以上版本(推荐)
  \end{itemize}
\item
  或其他支持运行CRESM的操作系统
\end{itemize}

\hypertarget{ux786cux4ef6ux914dux7f6eux8981ux6c42}{%
\subsection{硬件配置要求}\label{ux786cux4ef6ux914dux7f6eux8981ux6c42}}

推荐配置（具体需求取决于区域大小和数据规模）：

\begin{itemize}
\item
  CPU：48核及以上
\item
  内存：100 GB RAM
\item
  磁盘空间：至少 100 GB 可用空间
\end{itemize}

\hypertarget{ux7cfbux7edfux8f6fux4ef6ux4f9dux8d56}{%
\subsection{系统软件依赖}\label{ux7cfbux7edfux8f6fux4ef6ux4f9dux8d56}}

CPS 运行过程中依赖多个底层科学计算库和并行计算组件，请确保以下软件已正确安装：

\begin{itemize}
\item
  Intel Fortran 编译器套件
\item
  Intel MPI 并行环境
\item
  NetCDF-C 与 NetCDF-Fortran 库
\item
  HDF5 库
\item
  netCDF Operator (NCO)
\item
  Climate Data Operators (CDO)
\item
  NCAR Command Language (NCL)
\end{itemize}

建议用户在系统环境中验证以下工具可正常使用：

\begin{itemize}
\item
  \texttt{mpirun}
\item
  \texttt{ncdump}
\item
  \texttt{ncl}
\end{itemize}

\hypertarget{python-ux8fd0ux884cux73afux5883}{%
\subsection{Python 运行环境}\label{python-ux8fd0ux884cux73afux5883}}

CPS 基于 Python 运行，建议使用 Conda 进行环境管理，具体环境配置请参考5.1节。 Python 版本要求：

\begin{itemize}
\tightlist
\item
  Python 3.9 及以上版本
\end{itemize}

主要依赖库包括：

\begin{itemize}
\item
  numpy
\item
  pandas
\item
  matplotlib
\item
  cartopy
\item
  Pillow
\item
  geopandas
\item
  fiona
\item
  esmpy
\item
  esmf
\item
  xesmf
\item
  pyhdf
\item
  pyparsing
\item
  python-dateutil
\end{itemize}

\hypertarget{ux6a21ux5f0fux7cfbux7edfux4f9dux8d56}{%
\subsection{模式系统依赖}\label{ux6a21ux5f0fux7cfbux7edfux4f9dux8d56}}

在运行 CPS 之前，需要提前完成以下模式组件的安装与编译配置：

\begin{itemize}
\item
  区域气候模式 CRESM
\item
  区域气候模式CWRF的前处理套件(CWPS)
\end{itemize}

上述组件必须能够独立正常运行，以保证 CPS 前处理流程的正确执行。 我们提供一个标准的工具库CRESM\_ToolBox去完成上述系统依赖安装。

\hypertarget{ux8fd0ux884cux73afux5883ux914dux7f6eux8981ux6c42}{%
\subsection{运行环境配置要求}\label{ux8fd0ux884cux73afux5883ux914dux7f6eux8981ux6c42}}

推荐使用单独的环境配置文件（例如: \texttt{\textasciitilde{}/.bashrc\_cresm}）管理CPS运行环境。

\hypertarget{ux76eeux5f55ux7ed3ux6784}{%
\section{目录结构}\label{ux76eeux5f55ux7ed3ux6784}}

本节介绍 CPS 的主要目录及其用途。

\hypertarget{ux9879ux76eeux6839ux76eeux5f55ux7ed3ux6784}{%
\subsection{项目根目录结构}\label{ux9879ux76eeux6839ux76eeux5f55ux7ed3ux6784}}

CPS 项目目录结构如下：

\begin{snugshade}
\begin{CPSCode}
CPS/
    ├── PrepScript/          # 核心前处理脚本
    ├── SpinUpScript/        # CoLM Spin-Up脚本
    ├── CRESM_ToolBox/       # CRESM 工具箱
    ├── Case/                # 个例与实验工作目录
    └── README.md            # 使用说明文档
\end{CPSCode}
\end{snugshade}

\hypertarget{ux6838ux5fc3ux7a0bux5e8fux76eeux5f55prepscript}{%
\subsection{核心程序目录（PrepScript）}\label{ux6838ux5fc3ux7a0bux5e8fux76eeux5f55prepscript}}

\begin{snugshade}
\begin{CPSCode}
PrepScript/
    ├── CRESM_Preprocessing_System.py   # 主入口程序
    ├── case.ini                        # 实际实验配置
    ├── env.ini                         # 实际环境配置
    ├── Config/                         # 配置参考文件
    │   ├── case.ini.template           # 实验配置参考
    │   ├── env.ini.gnu_tms             # GNU TMS 编译环境参考
    │   ├── env.ini.gnu_hydro           # GNU Hydro 编译环境参考
    │   └── env.ini.intel               # Intel 编译环境参考
    ├── Modules/                        # Python 前处理模块
    │   ├── PrepCWRF.py                 # CWRF 前处理模块
    │   ├── PrepCoLM.py                 # CoLM 前处理模块
    │   ├── PrepCRESM.py                # CPL7 映射文件生成模块
    │   └── Utils/                      # 公共工具库
    │       ├── Consts.py
    │       ├── ICBC.py
    │       ├── Logger.py
    │       ├── Tools.py
    │       └── Domain.py                # 独立研究区域诊断模块
    └── Resources/                      # 模板、外部脚本与驱动信息
        ├── Forcing/
        ├── NML/
        └── ProcessScript/
\end{CPSCode}
\end{snugshade}

功能说明：

\begin{itemize}
\item
  \texttt{CRESM\_Preprocessing\_System.py} 系统主入口脚本，负责流程调度与模块调用。
\item
  \texttt{PrepCWRF.py} 用于执行 CWRF 静态数据处理、IC/BC 生成等流程。
\item
  \texttt{PrepCoLM.py} 负责陆面模式 CoLM 的基础数据准备与初始场构建。
\item
  \texttt{PrepCRESM.py} 用于生成耦合器所需的网格映射和插值权重文件。
\item
  \texttt{Domain.py} 用于基于 ERA5 等诊断数据计算异常场、相关系数和显著性研究区域。
\end{itemize}

\hypertarget{ux5b9eux9a8cux5de5ux4f5cux76eeux5f55case}{%
\subsection{实验工作目录（Case）}\label{ux5b9eux9a8cux5de5ux4f5cux76eeux5f55case}}

\begin{snugshade}
\begin{CPSCode}
Case/
├──[GridName]/                    # 实验数据汇总目录
    ├──Create_Run_From_CPS.py     # 用以创建运行案例的脚本
    ├──Grid_[GridName]/           # 静态数据
    └──ICBC_[GridName]/           # CWRF初始场及驱动场数据
├──PrepCWRF/                      # CWRF相关数据制作目录
    ├──[GridName]/                # CWRF相关数据汇总目录
    ├──First_StaticData/          # 静态数据制作目录
        └──*Geog_[GridName]/      # CWRF geog数据目录 (对同一区域可复用)
    └──Second_ICBC/               # CWRF初始场及驱动场数据制作目录
├──PrepCoLM/                      # CoLM相关数据制作目录
    ├──[GridName]/                # CoLM相关数据汇总目录
    ├──First_GenMesh/             # CoLM与CWRF网格映射文件制作目录
    ├──Second_MakeSrf/            # CoLM相关静态数据制作目录
        └──*CoLMSrf_[GridName]    # CWRF geog数据目录 (对同一区域和define配置可复用)
    └──Third_Remap/               # CoLM数据remap至CWRF网格
├──PrepCRESM/                     # CRESM相关数据制作目录
    └──[GridName]/                # cpl7相关数据汇总目录
├──Domain/                        # 独立研究区域诊断结果
    ├──ERA5_Anomalies_[GridName].nc
    ├──ERA5_Correlation_[GridName].nc
    ├──ERA5_Correlation_PValues_[GridName].nc
    ├──Correlation_Plots/            # 相关系数图
    └──Study_Area_Plots/             # 显著性研究区域图
├──NMLS/                          # namelist
└──Log/                           # 日志文件目录
\end{CPSCode}
\end{snugshade}

说明：

\begin{itemize}
\item
  每个 Case 目录对应一个独立实验区域或分辨率配置
\item
  \texttt{Geog\_{[}GridName{]}} 和\texttt{CoLMSrf\_{[}GridName{]}}分别为CWRF和CoLM的静态数据目录，在同一实验配置下可复用。 实验配置由唯一的Geog.ID验证。
\item
  系统会在该目录下自动生成各模块运行所需的中间文件和结果数据。
\item
  用户通常只需要管理 Case 层级目录，无需修改内部自动生成文件。
\item
  \textbf{警告:} 如果在实验配置文件中启用了自动清理缓存(CleanTempFiles)， 那只有 \texttt{Case/{[}GridName{]}/} 文件夹会被保留，请谨慎使用。
\end{itemize}

\hypertarget{cpsux6d41ux7a0bux603bux89c8}{%
\section{CPS流程总览}\label{cpsux6d41ux7a0bux603bux89c8}}

CPS提供了一个集成且自动化的工作流程，该工作流程旨在以统一且可复现的方式处理大气强迫、陆面数据集和耦合器映射文件的预处理。CPS 通过用户定义的配置文件控制的分阶段管道执行预处理任务，总体执行流程如下

\begin{snugshade}
\begin{CPSCode}
        case.ini + env.ini                # 用户配置文件
                 ↓
    ┌──────────────────────────┐
    │  CRESM Preprocessing     │          # 主程序
    │        System            │
    └──────────────────────────┘
                 ↓
    ┌──────────────────────────┐
    │  PrepCWRF                │          # 准备CWRF所需数据
    │  - Domain Visualization  │          # 可视化研究域
    │  - Static Geogrid        │          # geogrid
    │  - Vegetation Parameters │          # 植被参数(FVC, LAI, SAI) 
    │  - Initial/Boundary Data │          # 初始场/侧边界驱动场
    └──────────────────────────┘
                 ↓
    ┌──────────────────────────┐
    │  PrepCoLM                │          # 准备CoLM所需数据
    │  - Mesh Generation       │          # 制作映射文件
    │  - Surface Data          │          # 准备CoLM的陆面数据
    │  - Initial Conditions    │          # 初始化CoLM
    │  - Spin-up & Remapping   │          # Spin-up和重网格
    └──────────────────────────┘
                 ↓
    ┌──────────────────────────┐
    │  PrepCRESM               │          # 准备CRESM所需数据
    │  - Coupler Mapping Files │          # 耦合器权重文件
    └──────────────────────────┘
\end{CPSCode}
\end{snugshade}

\hypertarget{ux5b89ux88c5ux4e0eux90e8ux7f72}{%
\section{安装与部署}\label{ux5b89ux88c5ux4e0eux90e8ux7f72}}

\hypertarget{python-ux73afux5883ux90e8ux7f72}{%
\subsection{Python 环境部署}\label{python-ux73afux5883ux90e8ux7f72}}

CPS 的 Python 运行环境依赖多个科学计算、地理信息处理和数据可视化相关库。为保证依赖版本的一致性与跨平台兼容性，所有 Python 依赖已统一维护在CRESM\_ToolBox目录的三个环境配置文件中。为防止环境冲突，我们需要新建3个python环境。

\emph{请注意在现有版本中，用户只能通过conda进行环境配置。}

用户无需手动逐一安装各个依赖库，只需通过以下方式完成环境配置：

\begin{snugshade}
\begin{CPSCode}
conda env create -f cresm.yml              

conda env create -f cresm_xesmf.yml

conda env create -f Chaomodis.yml
\end{CPSCode}
\end{snugshade}

以上命令会新建三个新环境：

\begin{itemize}
\item
  \texttt{cresm}: 负责运行\texttt{main}, \texttt{FlowDir}, \texttt{plot}等脚本
\item
  \texttt{cresm\_xesmf}: 负责运行\texttt{GenMesh}, \texttt{Remap}等脚本
\item
  \texttt{Chaomodis}: 负责运行\texttt{FVC}, \texttt{IGBP}, \texttt{LAI/SAI}
\end{itemize}

\hypertarget{cresm-ux5de5ux5177ux7bb1ux5b89ux88c5}{%
\subsection{CRESM 工具箱安装}\label{cresm-ux5de5ux5177ux7bb1ux5b89ux88c5}}

CPS提供了一个完整的CRESM工具箱，该工具箱包含区域气候模式CWRF的前处理系统CWPS和依赖数据集，CWRF陆面数据制作相关代码及依赖数据。我们为用户提供了便捷的安装脚本\texttt{SetupAll.py}. 用户可以在设置完\texttt{Makeoptions.ini}之后通过\texttt{SetupAll.py}一次安装部署所有工具。\texttt{Makeoptions.ini}包含：

\begin{snugshade}
\begin{CPSCode}
[LIBPATHS]
NETCDF_path = $software/netcdf-single-r        # netcdf 安装目录   
JASPER_path = $software/jasper                 # jasper 安装目录
PNG_path    = $software/libpng                 # libpng 安装目录 
ZLIB_path   = $software/zlib-1.2.11            # zlib 安装目录
\end{CPSCode}
\end{snugshade}

配置完上述路径后，您可以通过以下命令完成安装：

\begin{snugshade}
\begin{CPSCode}
python SetupAll.py
\end{CPSCode}
\end{snugshade}

\hypertarget{ux914dux7f6eux6587ux4ef6}{%
\section{配置文件}\label{ux914dux7f6eux6587ux4ef6}}

CRESM Data Preparation System（CPS）运行时读取 \texttt{case.ini} 和 \texttt{env.ini}，分别用于描述具体实验流程和运行环境依赖。\texttt{Config/} 目录中的 \texttt{case.ini.template} 及 \texttt{env.ini.*} 仅作为参考文件，不是运行时直接读取的配置文件。

\hypertarget{ux7cfbux7edfux73afux5883ux914dux7f6eux6587ux4ef6-env.ini}{%
\subsection{\texorpdfstring{系统环境配置文件 （\texttt{env.ini}）}{系统环境配置文件 （env.ini）}}\label{ux7cfbux7edfux73afux5883ux914dux7f6eux6587ux4ef6-env.ini}}

环境配置文件用于描述``在哪里运行 + 用什么运行''，与具体实验无关。运行时 CPS 固定读取 \texttt{env.ini}。如需切换运行环境，请参考 \texttt{Config/} 中的环境参考文件，并将相应内容更新到 \texttt{env.ini}：

\begin{longtable}[]{@{}p{0.48\linewidth}p{0.48\linewidth}@{}}
\toprule
文件 & 适用环境 \\
\midrule
\endhead
\texttt{Config/env.ini.gnu\_tms} & GNU TMS 环境参考 \\
\texttt{Config/env.ini.gnu\_hydro} & GNU Hydro 环境参考 \\
\texttt{Config/env.ini.intel} & Intel 环境参考 \\
\bottomrule
\end{longtable}

选定的环境配置文件包含以下配置：

\hypertarget{environment-system-environment-information}{%
\subsubsection{\texorpdfstring{\textbf{{[}Environment{]}} -- System environment information}{{[}Environment{]} -- System environment information}}\label{environment-system-environment-information}}

\begin{longtable}[]{@{}p{0.23\linewidth}p{0.09\linewidth}p{0.43\linewidth}p{0.18\linewidth}@{}}
\toprule
参数名 & 类型 & 说明 & 示例 \\
\midrule
\endhead
SYS\_CWRF & str & CWRF 执行所需的环境路径（source 加载） & \textasciitilde/.cresm \\
SYS\_CoLM & str & CoLM 执行所需的环境路径（source 加载） & \textasciitilde/.CoLM\_2024 \\
SYS\_NCL & str & NCL 执行所需的环境路径（source 加载） & \textasciitilde/.cresm \\
SYS\_ESMF & str & ESMF 执行所需的环境路径（Use\_ESMF=True 时 source 加载） & \textasciitilde/.cresm \\
CONDA\_CRESM & str & CRESM 前处理所需的 Conda 环境名称 & cresm \\
CONDA\_XESMF & str & 5.1节安装的Conda 的\texttt{cresm\_xesmf}的环境名称 & cresm\_xesmf \\
CONDA\_CHAO & str & 5.1节安装的Conda 的\texttt{Chaomodis}的环境名称 & Chaomodis \\
CONDA\_UNGRIB & str & 5.1节安装的Conda 的\texttt{ungrib}的环境名称 & ungrib \\
\bottomrule
\end{longtable}

\hypertarget{paths-system-environment-information}{%
\subsubsection{\texorpdfstring{\textbf{{[}Paths{]}} -- System environment information}{{[}Paths{]} -- System environment information}}\label{paths-system-environment-information}}

\begin{longtable}[]{@{}p{0.23\linewidth}p{0.09\linewidth}p{0.42\linewidth}p{0.18\linewidth}@{}}
\toprule
参数名 & 类型 & 说明 & 示例 \\
\midrule
\endhead
ScriptPath & str & \texttt{PrepScript}文件夹的绝对路径 & CPS/PrepScript \\
CoLMModelPath & str & \texttt{CoLM}源代码/已编译模型根目录的路径 & CRESM/MODEL/CoLM/CoLM202X \\
CoLMRawDataPath & str & \texttt{CoLM}原始地理/土壤数据集的路径 & /shr03/CoLMrawdata \\
CoLMRunDataPath & str & \texttt{CoLM}运行时所需数据集的路径 & /shr03/CoLMruntime \\
CoLMForcingPath & str & \texttt{CoLM}强迫数据的路径（例如：ERA5-Land） & /shr03/CoLM\_Forcing/ERA5LAND/ \\
RootToolBox & str & \texttt{CRESM}工具箱的根目录 & CPS/CRESM\_ToolBox \\
CWPSPath & str & \texttt{CWRF-CWPS} 工具的路径 & \$\{Paths:RootToolBox\}/ToolBox/CWRF-CWPS \\
CWRFToolPath & str & \texttt{CWRF} 工具箱的路径 & \$\{Paths:RootToolBox\}/ToolBox/CWRF\_Tools \\
GeogDataPath & str & \texttt{CWRF\ Geogrid} 二进制数据的路径 & geog\_wm\_modified\_lake \\
LandSeaMaskPath & str & 陆地---海洋掩膜数据的路径 & land\_sea\_mask \\
CWPSStaticPath & str & \texttt{CWPS}静态表/文件的路径 & CWPS\_static \\
GlobalLakeDepth & str & 全球湖泊深度数据文件路径 & newGlobalLakeDepth.dat \\
GlobalLakeStatus & str & 全球湖泊状态数据文件的路径 & newGlobalLakeStatus.dat \\
WMEJUngrib & str & \texttt{WMEJ\_NC2IM} 工具的路径 & \$\{Paths:RootToolBox\}/ToolBox/WMEJ\_NC2IM \\
WMEJModis & str & \texttt{CoLM\ LAI} 数据处理工具的路径 & \$\{Paths:RootToolBox\}/ToolBox/WMEJ\_modis \\
ChaoModis & str & \texttt{MODIS}处理工具的路径 & \$\{Paths:RootToolBox\}/ToolBox/Chao\_modis \\
NCOPath & str & \texttt{ncks} 可执行文件的路径 & /usr/bin/ncks \\
CDOPath & str & \texttt{cdo} 可执行文件的路径 & /usr/bin/cdo \\
NCLPath & str & \texttt{ncl} 可执行文件的路径 & /usr/bin/ncl \\
ESMFWeightGenPath & str & \texttt{ESMF\_RegridWeightGen} 可执行文件的路径（Use\_ESMF=True 时使用） & /path/to/ESMF\_RegridWeightGen \\
\bottomrule
\end{longtable}

\hypertarget{forcingname-system-environment-information}{%
\subsubsection{\texorpdfstring{\textbf{{[}ForcingName{]}} -- System environment information}{{[}ForcingName{]} -- System environment information}}\label{forcingname-system-environment-information}}

\begin{longtable}[]{@{}p{0.23\linewidth}p{0.09\linewidth}p{0.43\linewidth}p{0.18\linewidth}@{}}
\toprule
参数名 & 类型 & 说明 & 示例 \\
\midrule
\endhead
ForcingDataName & str & 驱动名称 & ERA5 \\
DataCalender & str & 驱动数据使用的日历类型（如 standard、noleap） & standard \\
Forc\_Info & str & 驱动信息相关文件（只有NC格式数据需要，GRIB填None） & None \\
Forc\_2D\_Path & str & 2D平面场文件路径模板（需要完整路径） & \texttt{SFC\textless{}YYYY\textgreater{}\textless{}MM\textgreater{}\textless{}DD\textgreater{}.grib} \\
Forc\_3D\_Path & str & 3D高空场文件路径模板（需要完整路径） & \texttt{PRESS\textless{}YYYY\textgreater{}\textless{}MM\textgreater{}\textless{}DD\textgreater{}.grib} \\
Forc\_SST\_Path & str & 海表温度场文件路径模板（需要完整路径） & \texttt{SST\textless{}YYYY\textgreater{}\textless{}MM\textgreater{}\textless{}DD\textgreater{}.grib} \\
\bottomrule
\end{longtable}

\begin{quote}
注意： 用户可以同时定义多个强迫（例如，{[}cfsv2{]}、{[}gfs{]}）
\end{quote}

\hypertarget{domain-independent-study-domain-analysis}{%
\subsubsection{\texorpdfstring{\textbf{{[}Domain{]}} -- Independent Study-Domain Analysis}{{[}Domain{]} -- Independent Study-Domain Analysis}}\label{domain-independent-study-domain-analysis}}

\texttt{{[}Domain{]}} 用于 \texttt{-D/-\/-domain} 独立研究区域诊断。标准 CWRF-CoLM-CRESM 前处理流程不使用这些选项；启用诊断时，除研究区域边界参数外，以下选项均由主程序统一检查。

\begin{longtable}[]{@{}p{0.23\linewidth}p{0.09\linewidth}p{0.43\linewidth}p{0.18\linewidth}@{}}
\toprule
参数名 & 类型 & 说明 & 示例 \\
\midrule
\endhead
DefaultDataName & str & 诊断数据名称，用于结果文件命名 & ERA5 \\
DomainDataPath & file & 诊断分析使用的 NetCDF 数据文件 & \texttt{/path/to/era5.nc} \\
VarList & list & 参与分析的变量，使用逗号分隔 & \texttt{u,\ v,\ r,\ t,\ q} \\
StudyAreaCriteria & list & 定义研究区域的变量与气压层，格式为 \texttt{变量:气压层} & \texttt{u:850,\ v:850,\ r:925,\ q:925} \\
BoundaryFracThres & float & 显著区域覆盖比例阈值，范围为 \texttt{0.0-1.0} & \texttt{0.15} \\
LevelDimName & str & 气压层坐标名称 & \texttt{pressure\_level} \\
TimeDimName & str & 时间坐标名称 & \texttt{valid\_time} \\
LatDimName & str & 纬度坐标名称 & \texttt{latitude} \\
LonDimName & str & 经度坐标名称 & \texttt{longitude} \\
AnalyzeLonMin/Max & float & 诊断分析的经度范围 & \texttt{60\ /\ 140} \\
AnalyzeLatMin/Max & float & 诊断分析的纬度范围 & \texttt{0\ /\ 60} \\
\bottomrule
\end{longtable}

研究区域可以通过 Shapefile 或经纬度矩形框指定，命令行参数见第 7 节。

\hypertarget{ux5b9eux9a8cux914dux7f6eux6587ux4ef6-case.ini}{%
\subsection{\texorpdfstring{实验配置文件 （\texttt{case.ini}）}{实验配置文件 （case.ini）}}\label{ux5b9eux9a8cux914dux7f6eux6587ux4ef6-case.ini}}

\texttt{case.ini} 用于控制 ``做什么 + 在哪里 + 做多久''，决定 CPS 的执行流程与目标实验。用户可以参考 \texttt{Config/case.ini.template} 填写案例配置。

\hypertarget{baseinfo-global-workflow-settings}{%
\subsubsection{\texorpdfstring{\textbf{{[}BaseInfo{]}} -- Global Workflow Settings}{{[}BaseInfo{]} -- Global Workflow Settings}}\label{baseinfo-global-workflow-settings}}

\begin{longtable}[]{@{}p{0.23\linewidth}p{0.09\linewidth}p{0.42\linewidth}p{0.18\linewidth}@{}}
\toprule
参数名 & 类型 & 说明 & 示例 \\
\midrule
\endhead
CleanTempFiles & bool & 若为 True，则删除中间文件以节省磁盘空间。 & True/False \\
Use\_ESMF & bool & 若为 True，使用 NCL 生成 CPL7 Step1 描述文件，再由 ESMF 生成权重；若为 False，使用 NCL 全流程。 & True/False \\
Use\_CoLMLAI & bool & 若为 True，则使用 CoLM 的 LAI 数据而不是默认的 MODIS 数据制作 LAI。 & True/False \\
Use\_CoLMSeaMask & bool & 若为 True，则使用 CoLM 的海陆掩膜；若为 False，则使用 shpfile 湖泊/海洋掩膜。 & True/False \\
Enable\_TimeChunk & bool & 启用时间分块，对较长时间段的 ICBC 处理进行分块并行。 & True/False \\
TimeChunkCount & int & 时间分块总数。 & 4 \\
TimeChunkGroupSize & int & 每个串行分组中允许同时处理的最大时间块数量。 & 2 \\
Reuse\_Metgrid & bool & 若为 True，仅复用完成标记和全部 \texttt{met\_em} 文件均存在的时间块。 & True/False \\
Clean\_Ungrib & bool & 若为 True，在 Metgrid 和 SST 处理成功后删除 2D、3D、SST Ungrib 中间文件。 & True/False \\
\bottomrule
\end{longtable}

\texttt{{[}BaseInfo{]}} 中的所有选项均为必填项。\texttt{TimeChunkGroupSize} 表示每组的时间块数量， 各组依次执行，组内时间块并行执行。

当 \texttt{Reuse\_Metgrid\ =\ True} 时，CPS 会检查时间块目录中的 \texttt{.metgrid.complete} 和全部预期 \texttt{met\_em.d01.*.nc} 文件。检查通过后跳过该时间块的 CWPS 文件链接、Ungrib 和 Metgrid， 但后续的 Real 步骤仍会执行。

\texttt{PrepCRESM} 的 Step1 在 \texttt{Use\_ESMF=True} 时先由 NCL 生成网格描述文件，再由 \texttt{ESMF\_RegridWeightGen} 生成权重文件；Step2 继续使用 Python 完成 CoLM 索引重排。 将 \texttt{Use\_ESMF} 设置为 \texttt{False} 时，继续使用原有的 \texttt{generate\_cpl7\_wgt\_ALO\_step1.ncl} 完成描述文件和权重文件生成。

\hypertarget{prepcwrf-cwrf-preprocessing-switches}{%
\subsubsection{\texorpdfstring{\textbf{{[}PrepCWRF{]}} -- CWRF Preprocessing Switches}{{[}PrepCWRF{]} -- CWRF Preprocessing Switches}}\label{prepcwrf-cwrf-preprocessing-switches}}

\begin{longtable}[]{@{}p{0.23\linewidth}p{0.09\linewidth}p{0.43\linewidth}p{0.18\linewidth}@{}}
\toprule
参数名 & 类型 & 说明 & 示例 \\
\midrule
\endhead
CWRFCoreNum & int & 用于 CWRF MPI 任务的 CPU 核心数（建议 4-24 个）。 & 16 \\
Go\_ShowDomain & bool & 绘制域预览图（检查域覆盖率）。 & True/False \\
Go\_Geogrid & bool & 运行地理网格，生成静态地理数据。 & True/False \\
Go\_FVC & bool & 处理植被覆盖度数据。 & True/False \\
Go\_LAI & bool & 处理叶面积指数数据。 & True/False \\
Go\_IGBP & bool & 处理 IGBP 土地利用数据。 & True/False \\
Go\_SAI & bool & 处理茎面积指数数据。 & True/False \\
Collect\_GeogData & bool & 收集 CWRF 静态数据以方便后续复用。 & True/False \\
Go\_Ungrib & bool & 运行 Ungrib，解码 GRIB/NC 强迫数据。 & True/False \\
Go\_Metgrid & bool & 运行 Metgrid，将气象数据插值到模型网格。 & True/False \\
Go\_Real & bool & 运行 Real，生成 \texttt{wrfinput}/\texttt{wrfbdy}。 & True/False \\
Go\_VBS & bool & 运行 VBS，进行植被/反照率处理。 & True/False \\
Copy\_CWRF\_Output & bool & 将最终 CWRF 输入复制到案例目录。 & True/False \\
\bottomrule
\end{longtable}

\texttt{{[}PrepCWRF{]}} 中的所有选项均为必填项。

\hypertarget{prepcolm-colm-preprocessing-switches}{%
\subsubsection{\texorpdfstring{\textbf{{[}PrepCoLM{]}} -- CoLM Preprocessing Switches}{{[}PrepCoLM{]} -- CoLM Preprocessing Switches}}\label{prepcolm-colm-preprocessing-switches}}

\begin{longtable}[]{@{}p{0.23\linewidth}p{0.09\linewidth}p{0.42\linewidth}p{0.18\linewidth}@{}}
\toprule
参数名 & 类型 & 说明 & 示例 \\
\midrule
\endhead
CoLMCoreNum & int & 用于 CoLM MPI 任务的 CPU 核心数（建议 24-128 个） & 96 \\
Go\_MeshGrid & bool & 生成CoLM与CWRF的非结构化网格映射文件 & True/False \\
Go\_MakeSrf & bool & 生成表面数据集（mksrfdata） & True/False \\
Go\_MakeIni & bool & 生成初始条件（mkinidata） & True/False \\
Go\_CoLMTempRun & bool & 进行短时运行，生成history & True/False \\
Go\_Remap & bool & 重新映射history和restart文件 & True/False \\
Copy\_CoLM\_Output & bool & 将最终的 CoLM 输入复制到案例目录 （GridName/PrepCoLM/GridName） & True/False \\
\bottomrule
\end{longtable}

\hypertarget{prepcresm-coupler-weight-file}{%
\subsubsection{\texorpdfstring{\textbf{{[}PrepCRESM{]}} -- Coupler weight file}{{[}PrepCRESM{]} -- Coupler weight file}}\label{prepcresm-coupler-weight-file}}

\begin{longtable}[]{@{}p{0.23\linewidth}p{0.09\linewidth}p{0.43\linewidth}p{0.18\linewidth}@{}}
\toprule
参数名 & 类型 & 说明 & 示例 \\
\midrule
\endhead
Go\_Coupler\_Prep & bool & 为耦合器生成权重和映射 & True/False \\
\bottomrule
\end{longtable}

\texttt{{[}PrepCRESM{]}} 中的所有选项均为必填项。

\hypertarget{gatherdata-final-output-collection}{%
\subsubsection{\texorpdfstring{\textbf{{[}GatherData{]}} -- Final Output Collection}{{[}GatherData{]} -- Final Output Collection}}\label{gatherdata-final-output-collection}}

\begin{longtable}[]{@{}p{0.23\linewidth}p{0.09\linewidth}p{0.43\linewidth}p{0.18\linewidth}@{}}
\toprule
参数名 & 类型 & 说明 & 示例 \\
\midrule
\endhead
Collect\_CWRF\_Output & bool & 收集最终的CWRF输入数据 & True/False \\
Collect\_CoLM\_Output & bool & 收集最终的CoLM输入数据 & True/False \\
Collect\_CRESM\_Output & bool & 收集耦合器名称列表和映射文件（GridName/GridName） & True/False \\
\bottomrule
\end{longtable}

\texttt{{[}GatherData{]}} 中的所有选项均为必填项。

\hypertarget{gridname-domain-experiment-specifics}{%
\subsubsection{\texorpdfstring{\textbf{{[}GridName{]}} -- Domain \& Experiment Specifics}{{[}GridName{]} -- Domain \& Experiment Specifics}}\label{gridname-domain-experiment-specifics}}

每个实验 \texttt{{[}GridName{]}} 小节中，除 \texttt{LakeThreshold} 外的所有选项均必须存在。\texttt{CoLMNMLPath} 可以填写 \texttt{default} 或留空，表示使用内置模板。 \texttt{LakeThreshold} 可以省略或留空，系统将默认使用 \texttt{0.5}；如果填写，则取值范围必须为 \texttt{0.0-1.0}。

\texttt{CoLMNMLPath} 用于指定一个完整的 CoLM namelist 文件。CPS 会以该文件为基础，仅更新案例名称、研究区域边界、时间、路径、网格文件和输出频率等 CPS 管理的参数，用户在该文件中设置的其他 CoLM 参数会原样保留。\texttt{default} 或空值表示使用 \texttt{Resources/NML/unstructured\_cwrf.colm.ctl}。

\begin{longtable}[]{@{}p{0.23\linewidth}p{0.09\linewidth}p{0.42\linewidth}p{0.18\linewidth}@{}}
\toprule
参数名 & 类型 & 说明 & 示例 \\
\midrule
\endhead
CaseOutputPath & str & 案例输出文件的根目录 & /path/to/CPS/Case \\
CoLMNMLPath & path/default & 完整 CoLM namelist 文件的绝对路径；填写 \texttt{default} 或留空时使用内置模板。CPS 只替换自身负责的配置，其余配置保持不变。 & \texttt{default} \\
ForcingDataName & str & 驱动名称，必须与 \texttt{env.ini} 中的 forcing 小节对应（如 CFSV2、ERA5、MPI-ESM1-2-HR\_ssp245） & ERA5 \\
StartTime & str & 开始时间，格式为 YYYY-MM-DD\_HH:MM:SS & 2000-01-01\_00:00:00 \\
EndTime & str & 结束时间，格式为 YYYY-MM-DD\_HH:MM:SS & 2000-01-03\_00:00:00 \\
EdgeNum\_WE & int & 东西向网格数+1 & 232 \\
EdgeNum\_SN & int & 南北向网格数+1 & 172 \\
dx\_WE & float & 东西向网格分辨率 (单位: m) & 30000 \\
dy\_SN & float & 南北向网格分辨率 (单位: m) & 30000 \\
RefLat & float & 该区域的中心纬度 & 35.1778 \\
RefLon & float & 该区域的中心经度 & 110.0 \\
True\_Lat1 & float & 兰伯特投影的下切线纬度 & 30.0 \\
True\_Lat2 & float & 兰伯特投影的上切线纬度 & 60.0 \\
StandLon & float & 兰伯特投影的标准经度 & 110.0 \\
BdyWidth & int & 侧向边界松弛宽度，必须为不小于 13 的奇数。 & 15 \\
LakeThreshold & float & 可选湖泊比例（0.0-1.0），湖泊比例大于此值的网格单元将被视为湖泊点；缺失或留空时默认为 0.5。 & 0.5 \\
MeshSize & int & CoLM 与 CWRF 的非结构化网格映射文件分辨率，取值为 1、2 或 3。1=Coarse，2=Medium，3=Fine。 & 1 \\
\bottomrule
\end{longtable}

\begin{quote}
注意： 用户可以同时定义多个实验配置（例如，{[}CN\_15km{]}、{[}Yangtze{]}）
\end{quote}

上述 \texttt{case.ini} 的 Section 和选项检查均由 \texttt{CRESM\_Preprocessing\_System.py} 集中完成， 后续处理模块使用已经完成检查的配置，不再重复进行配置检查。

\hypertarget{ux57faux672cux7528ux6cd5}{%
\section{基本用法}\label{ux57faux672cux7528ux6cd5}}

\hypertarget{ux4f7fux7528ux524dux51c6ux5907}{%
\subsection{使用前准备}\label{ux4f7fux7528ux524dux51c6ux5907}}

在运行 CRESM 数据前处理系统之前，请完成以下准备工作：

\begin{enumerate}
\item 系统环境已正确配置
  \begin{itemize}
  \item 已完成 \texttt{env.ini} 中所有路径与环境的配置；
  \item CWRF、CoLM 和 Conda 环境均可正常激活；
  \item 外部程序（如 ncl、cdo、ncks）可在命令行调用。
  \end{itemize}
\item
  实验配置文件已准备

  \begin{itemize}
  \item
    参考 \texttt{Config/case.ini.template}，正确填写 \texttt{case.ini} 中的具体实验案例；
  \item
    确保 ForcingDataName 与 \texttt{env.ini} 中对应的 forcing 配置一致；
  \item
    检查模拟时间、网格参数、分辨率等关键信息。
  \end{itemize}
\end{enumerate}

\hypertarget{ux547dux4ee4ux884cux53c2ux6570ux8bf4ux660e}{%
\subsection{命令行参数说明}\label{ux547dux4ee4ux884cux53c2ux6570ux8bf4ux660e}}

CRESM 数据前处理系统通过统一的主程序启动，提供了诸多选项，用户可以使用\texttt{-h}选项查看命令行帮助

\textbf{命令}

\begin{snugshade}
\begin{CPSCode}
python CRESM_Preprocessing_System.py -h
\end{CPSCode}
\end{snugshade}

\textbf{输出(示例)}

\begin{snugshade}
\begin{CPSCode}
usage: CRESM_Preprocessing_System.py [-h] [-v] [-d] [-ch] [-e NAME] [-l] [-n NAME] [-g DIR]
                                     [-s DIR] [-y YEAR] [-c NAME] [-D]
                                     [-ds SHP] [-db LON_MIN LON_MAX LAT_MIN LAT_MAX]
                                     

CRESM Preprocessing System (CPS)

optional arguments:

  -h, --help        show this help message and exit
  -v, --version     Show version information and exit
  -d, --debug       Enable debug mode
  -ch, --confighelp Display configuration help and exit
  -l, --listcases   List available “useful cases” in case.ini and exit
  -n NAME, --gridname NAME
                    Specify the case name (matches a section in case.ini)
  -g DIR, --geogdir DIR
                    Use existing geography data from DIR for IC/BC preparation;
                    if omitted, new geography data will be generated
  -s DIR, --colmsrf DIR
                    Use existing CoLM surface data from DIR for PrepCoLM preparation;
                    if omitted, new surface data will be generated
  -y YEAR, --year YEAR  Override start/end year in the case (for yearly forcing data),
                    e.g., --year 2023
  -c NAME, --collectcase NAME
                    When using --year, gather all annual forcing data into the folder NAME;
                    inside that folder, each forcing file is renamed with the year suffix.
  -D, --domain       Run independent study-domain diagnostic analysis
  -ds SHP, -DS SHP, --domain-shp SHP
                    Use a Shapefile to define the diagnostic study area
  -db LON_MIN LON_MAX LAT_MIN LAT_MAX, -DB LON_MIN LON_MAX LAT_MIN LAT_MAX, --domain-bbox ...
                    Use a longitude/latitude bounding box to define the diagnostic study area
\end{CPSCode}
\end{snugshade}

\hypertarget{ux914dux7f6eux6587ux4ef6-1}{%
\subsubsection{配置文件}\label{ux914dux7f6eux6587ux4ef6-1}}

运行时 CPS 固定读取 \texttt{case.ini} 和 \texttt{env.ini}。如需切换 GNU TMS、GNU Hydro 或 Intel 环境， 请参考 \texttt{Config/} 中对应的环境参考文件，并将相应内容更新到 \texttt{env.ini}。

\hypertarget{v---version}{%
\subsubsection{\texorpdfstring{\texttt{-v,\ -\/-version}}{-v, -\/-version}}\label{v---version}}

\textbf{显示程序版本信息并退出。}

\begin{itemize}
\item
  用于确认当前使用的 CPS 版本
\item
  便于问题反馈与版本对齐
\end{itemize}

\hypertarget{d---debug}{%
\subsubsection{\texorpdfstring{\texttt{-d,\ -\/-debug}}{-d, -\/-debug}}\label{d---debug}}

\textbf{启用调试模式（Debug Mode）。}

\begin{itemize}
\item
  打印更详细的日志信息
\item
  包括路径检查、环境检查、每一步调用的外部命令
\item
  建议在出现错误时开启
\end{itemize}

\hypertarget{ch---confighelp}{%
\subsubsection{\texorpdfstring{\texttt{-ch,\ -\/-confighelp}}{-ch, -\/-confighelp}}\label{ch---confighelp}}

\textbf{显示配置文件说明并退出。}

\begin{itemize}
\item
  打印 \texttt{case.ini} 和 \texttt{env.ini} 的完整配置说明
\item
  用于快速查阅配置含义，无需打开手册
\end{itemize}

\hypertarget{l---listcases}{%
\subsubsection{\texorpdfstring{\texttt{-l,\ -\/-listcases}}{-l, -\/-listcases}}\label{l---listcases}}

\textbf{列出 \texttt{case.ini} 中已定义的可用实验案例并退出。}

\begin{itemize}
\item
  自动扫描 \texttt{case.ini}
\item
  显示所有可直接运行的 {[}GridName{]} 小节
\item
  不执行任何计算流程
\end{itemize}

\hypertarget{n-name---gridname-name}{%
\subsubsection{\texorpdfstring{\texttt{-n\ NAME,\ -\/-gridname\ NAME}}{-n NAME, -\/-gridname NAME}}\label{n-name---gridname-name}}

\textbf{指定要运行的实验案例名称。}

\begin{itemize}
\item
  NAME 必须与 \texttt{case.ini} 中的某个 {[}GridName{]} 小节一致
\item
  这是实际运行 CPS 时的核心参数
\item
\end{itemize}

\hypertarget{g-dir---geogdir-dir}{%
\subsubsection{\texorpdfstring{\texttt{-g\ DIR,\ -\/-geogdir\ DIR}}{-g DIR, -\/-geogdir DIR}}\label{g-dir---geogdir-dir}}

\textbf{使用已有的 CWRF 地理静态数据目录。} 若指定该参数：

\begin{itemize}
\item
  跳过 \texttt{Geogrid/FVC/LAI/SAI} 等静态地理数据生成步骤
\item
  直接使用 DIR 中已有的数据
\item
  网格配置有\texttt{Geog\_{[}GridName{]}}中的\texttt{Geog.ID}文件确认。不同网格配置不能混用。
\end{itemize}

Geog\_ID样式：

\begin{snugshade}
\begin{CPSCode}
X[232]_Y[172]_dx[30000.000000]_dy[30000.000000]_lat[35.177800]_lon[110.000000]_
tl1[30.000000]_tl2[60.000000]_slon[110.000000]_bw[15]_lk[0.500000]
\end{CPSCode}
\end{snugshade}

\hypertarget{s-dir---colmsrf-dir}{%
\subsubsection{\texorpdfstring{\texttt{-s\ DIR,\ -\/-colmsrf\ DIR}}{-s DIR, -\/-colmsrf DIR}}\label{s-dir---colmsrf-dir}}

\textbf{使用已有的 CoLM 地表数据目录。} 若指定该参数：

\begin{itemize}
\item
  跳过 \texttt{PrepCoLM} 中的地表数据生成流程\texttt{MakeSrf}
\item
  网格配置有\texttt{CoLMSrf\_{[}GridName{]}}中的\texttt{Geog.ID}文件确认。不同网格配置不能混用。
\item
  多案例共享同一套 CoLM surface 数据建议使用
\end{itemize}

\hypertarget{y-year---year-year}{%
\subsubsection{\texorpdfstring{\texttt{-y\ YEAR,\ -\/-year\ YEAR}}{-y YEAR, -\/-year YEAR}}\label{y-year---year-year}}

\textbf{覆盖\texttt{case.ini}中定义的起止年份，仅运行指定年份， 这会自动生成新实验\texttt{{[}GridName{]}.year}。}

\begin{itemize}
\item
  YEAR 为四位整数（如 2023）
\item
  自动调整：

  \begin{itemize}
  \item
    StartTime (如：2023-12-01\_00:00:00)
  \item
    EndTime (如：2025-01-05\_00:00:00)
  \item
    起止时间为\texttt{year-1年12月1日}至\texttt{year+1年1月5日}
  \end{itemize}
\item
  主要用于逐年强迫数据处理, 或跨年运行。
\end{itemize}

\hypertarget{c-name---collectcase-name}{%
\subsubsection{\texorpdfstring{\texttt{-c\ NAME,\ -\/-collectcase\ NAME}}{-c NAME, -\/-collectcase NAME}}\label{c-name---collectcase-name}}

\textbf{年度模式下的数据汇总参数（与 \texttt{-\/-year} 联合使用）。}

\begin{itemize}
\item
  将多个年度的强迫数据汇总到同一目录 NAME
\item
  每个文件自动追加年份后缀
\item
  便于后续多年模拟或批处理 \textbf{注意：仅在使用 \texttt{-\/-year} 后单独使用才有效}
\end{itemize}

\hypertarget{d---domain}{%
\subsubsection{\texorpdfstring{\texttt{-D,\ -\/-domain}}{-D, -\/-domain}}\label{d---domain}}

\textbf{运行独立研究区域诊断分析。}

该选项调用 \texttt{Domain.py}，使用 \texttt{env.ini} 的 \texttt{{[}Domain{]}} 配置读取诊断数据，依次生成月异常场、 研究区域相关系数、相关系数图以及显著性研究区域图。它不执行常规的 CWRF、CoLM 和 CPL7 前处理步骤，也不改变已有的模式输入文件。

研究区域必须通过以下两种方式之一指定：

\begin{itemize}
\item
  \texttt{-ds\ SHP} 或 \texttt{-DS\ SHP}：使用 Shapefile。Shapefile 需要包含有效的坐标参考系，系统会转换为 WGS84。
\item
  \texttt{-db\ LON\_MIN\ LON\_MAX\ LAT\_MIN\ LAT\_MAX} 或大写别名 \texttt{-DB}：使用经纬度矩形框。
\end{itemize}

示例：

\begin{snugshade}
\begin{CPSCode}
python CRESM_Preprocessing_System.py -n Yangtze_C_6km_MPI -D -db 90 120 20 35
python CRESM_Preprocessing_System.py -n Yangtze_C_6km_MPI -D -ds /path/to/study_area.shp
\end{CPSCode}
\end{snugshade}

两种边界同时提供时，Shapefile 优先。诊断结果保存在案例目录的 \texttt{Case/{[}GridName{]}/Domain/} 下；其中包括异常场和相关系数 NetCDF 文件，以及 \texttt{Correlation\_Plots/}、\texttt{Study\_Area\_Plots/} 两类图像目录。

\hypertarget{ux8fd0ux884c}{%
\subsection{运行}\label{ux8fd0ux884c}}

\hypertarget{ux6700ux5c0fux8fd0ux884cux793aux4f8b}{%
\subsubsection{最小运行示例}\label{ux6700ux5c0fux8fd0ux884cux793aux4f8b}}

当第一次运行一个新案例时推荐使用以下指令运行：

\begin{snugshade}
\begin{CPSCode}
python CRESM_Preprocessing_System.py -n CN_30km
\end{CPSCode}
\end{snugshade}

\hypertarget{ux5236ux4f5cux591aux5e74ux6570ux636e}{%
\subsubsection{制作多年数据}\label{ux5236ux4f5cux591aux5e74ux6570ux636e}}

当需要制作多年数据时，推荐先只运行5天，生成静态数据\texttt{CoLMSrf\_{[}GridName{]}} 和 \texttt{Geog\_{[}GridName{]}}。

\begin{snugshade}
\begin{CPSCode}
python CRESM_Preprocessing_System.py -n CN_30km
\end{CPSCode}
\end{snugshade}

之后复用静态数据再生成多年数据 (以2000年为例)

\begin{snugshade}
\begin{CPSCode}
python CRESM_Preprocessing_System.py -n CN_30km -y 2000 \
                                     -g path/to/Geog_[GridName] \
                                     -s path/to/CoLMSrf_[GridName]
\end{CPSCode}
\end{snugshade}

所有年份运行完毕后，再使用收集模式对所有年度数据进行收集

\begin{snugshade}
\begin{CPSCode}
python CRESM_Preprocessing_System.py -c CN_30km
\end{CPSCode}
\end{snugshade}

\hypertarget{ux8c03ux8bd5ux8fd0ux884c}{%
\subsubsection{调试运行}\label{ux8c03ux8bd5ux8fd0ux884c}}

出现错误时，可开启调试模式定位失败的步骤。建议仅开启需要排查的流程(\texttt{Go\_}) ，并启用\texttt{-d}模式

\begin{snugshade}
\begin{CPSCode}
python CRESM_Preprocessing_System.py -n CN_30km -d
\end{CPSCode}
\end{snugshade}

\hypertarget{ux521bux5efaux8fd0ux884cux6848ux4f8b}{%
\section{创建运行案例}\label{ux521bux5efaux8fd0ux884cux6848ux4f8b}}

通过CPS制作的数据既可以使用以前的\texttt{create\_run.py}创建运行案例，也可以使用更方便的\texttt{Create\_Run\_From\_CPS.py}创建运行案例。\texttt{Create\_Run\_From\_CPS.py}有以下可选参数：

\begin{snugshade}
\begin{CPSCode}
Create run directory for CRESM

optional arguments:
  -h, --help            show this help message and exit
  -casename CASENAME, --casename CASENAME
                        name for entire run
  -casedir CASEDIR, --casedir CASEDIR
                        The parent directory where the case will be created
  -griddir GRIDDIR, --griddir GRIDDIR
                        Full path to Grid directory
  -icbcdir ICBCDIR, --icbcdir ICBCDIR
                        Full path to ICBC directory
  -st STARTTIME, --starttime STARTTIME
                        YYYYMMDDHH
  -et ENDTIME, --endtime ENDTIME
                        YYYYMMDDHH
  -ncpus NCPUS, --ncpus NCPUS
                        number of cpus
  -jms {lsf,slurm}, --job_management_system {lsf,slurm}
  -ini INI, --ini INI   config.ini path
  -yearly, --yearly     If true, link ICBC files without renaming 
                        and copy wrfinput for start year
  -fvcom FVCOM, --fvcom FVCOM
                        True/False
\end{CPSCode}
\end{snugshade}

其中，\texttt{casename}是案例名称，\texttt{casedir}是案例创建目录，\texttt{griddir}是静态数据目录(即\texttt{Grid\_{[}GridName{]}})。\texttt{icbcdir}是CWRF初始化和侧边界场数据。\texttt{starttime}和\texttt{endtime}分别是开始时间和结束时间。\texttt{ncpus}是使用的核心数。\texttt{job\_management\_system}是任务管理系统，目前仅支持\texttt{lsf}, \texttt{slurm}, \texttt{ini} 是prep.ini文件的路径。\texttt{yearly}用以区分是否是分年的ICBC。如果数据是分年(即有年份后缀)，但是未启用，则只链接第一年的数据。

\hypertarget{ux65e5ux5fd7}{%
\section{日志}\label{ux65e5ux5fd7}}

CPS 分别记录主日志和各处理步骤的流程日志，用于检查运行状态和排查错误。 日志系统主要由\textbf{主日志（Main Log）和流程日志（Process Logs）}两部分组成。

\hypertarget{ux4e3bux65e5ux5fd7main-log}{%
\subsection{主日志（Main Log）}\label{ux4e3bux65e5ux5fd7main-log}}

\textbf{日志位置} 主日志文件默认保存在当前运行目录，文件命名格式为：

\begin{snugshade}
\begin{CPSCode}
DataPrepare.[GridName].log
\end{CPSCode}
\end{snugshade}

其中 \texttt{GridName} 对应命令行参数 \texttt{-\/-gridname} 指定的案例名称。 \textbf{主日志用于记录整个 CPS 工作流的全局信息，主要包括：}

\begin{itemize}
\item
  程序启动与结束信息
\item
  配置文件读取与参数解析结果
\item
  系统环境与路径检查结果
\item
  各预处理模块的开始与结束状态
\item
  外部程序调用命令（Debug 模式下）
\item
  全局错误与异常回溯信息 \textbf{日志级别:}
\item
  \texttt{INFO}：关键流程节点与总体状态信息
\item
  \texttt{DEBUG}：详细执行信息（仅在启用 --debug 时输出）
\item
  \texttt{WARNING}：非致命问题提示
\item
  \texttt{ERROR}：流程中断或执行失败信息
\end{itemize}

\textbf{建议：}

\begin{itemize}
\item
  日常生产运行可关闭 Debug
\item
  初次部署或出错时开启 \texttt{-\/-debug}
\end{itemize}

\hypertarget{ux6d41ux7a0bux65e5ux5fd7process-logs}{%
\subsection{流程日志（Process Logs）}\label{ux6d41ux7a0bux65e5ux5fd7process-logs}}

\textbf{日志位置} 主日志文件默认保存在当前运行目录，文件命名格式为：

\begin{snugshade}
\begin{CPSCode}
CaseOutputPath/[GridName]/Log/
\end{CPSCode}
\end{snugshade}

该目录用于集中存放各个处理步骤的标准输出与错误输出。

\textbf{命名规则} 流程日志采用如下统一命名格式：

\begin{snugshade}
\begin{CPSCode}
log.<ProcessName>[.<TimeTag>]
\end{CPSCode}
\end{snugshade}

在涉及时间分片或多次调用的流程中，日志文件可能附带时间标识或编号后缀。

\textbf{日志内容} 流程日志主要记录：

\begin{itemize}
\item
  外部可执行程序（如 \texttt{geogrid.exe}, \texttt{ungrib.exe}）的标准输出
\item
  MPI 并行运行信息
\item
  模型或工具自身产生的错误信息
\item
  详细的数值或文件处理过程 这些日志通常是定位失败原因的第一手信息。
\end{itemize}

\hypertarget{ux65e5ux5fd7ux4e0eux9519ux8befux5b9aux4f4dux5efaux8bae}{%
\subsection{日志与错误定位建议}\label{ux65e5ux5fd7ux4e0eux9519ux8befux5b9aux4f4dux5efaux8bae}}

当 CPS 执行失败时，建议按以下顺序排查：

\begin{enumerate}
\item 查看主日志
  \begin{itemize}
  \item 确认失败的流程阶段及路径（参见 \texttt{Error at:} 提示）；
  \item 检查环境、路径或参数相关错误。
  \end{itemize}
\item
  查看对应流程日志

  \begin{itemize}
  \item
    打开 \texttt{CaseOutputPath/{[}GridName{]}/Log/log.\textless{}ProcessName\textgreater{}}
  \item
    检查是否存在：

    \begin{itemize}
    \item
      库缺失（如 libnetcdf*.so）
    \item
      输入数据不存在
    \item
      MPI 或并行错误
    \end{itemize}
  \end{itemize}
\item
  开启 Debug 模式重试。若启用 \texttt{CleanTempFiles = True}，部分中间文件可能被自动删除；调试时请关闭该选项。
\end{enumerate}

\hypertarget{ux5e38ux89c1ux9519ux8befux6392ux67e5ux5efaux8bae}{%
\subsection{常见错误排查建议}\label{ux5e38ux89c1ux9519ux8befux6392ux67e5ux5efaux8bae}}

在目前版本中，部分操作系统可能会出现以下错误： \textbf{1. 无法找到依赖库，如netcdf}

\begin{itemize}
\tightlist
\item
  请确保在运行CPS前，source对应环境后重试，目前在CPS内部切换环境的功能并不稳定
\end{itemize}

\textbf{2. CWRF制作ICBC数据时出错且原因未知}

\begin{itemize}
\tightlist
\item
  部分网格的划分数会与分配的核心数冲突，因此可以尝试更改分配的核心数
\end{itemize}

\textbf{3. CoLM 部分运行失败且原因未知}

\begin{itemize}
\tightlist
\item
  可能是环境加载有误，请重新加载环境后再次运行
\end{itemize}

\begin{quote}
注意： 请注意，CoLM的define选项会影响你的目标实验配置，请确保define正确
\end{quote}
\end{document}
